\chapter{Superficies en el espacio}

\begin{definicion}
    Una superficie en $\mathbb{R}^3$ es un subconjunto $S\subset\mathbb{R}^3$ no vacío tal que $\forall p\in S$ existen:
    \begin{enumerate}[label=\alph*)]
        \item $U\subset\mathbb{R}^2$ abierto.
        \item Un entorno abierto $V$ de $p$ en $S$.
        \item Una aplicación $X:U\to \mathbb{R}^3$ diferenciable.
    \end{enumerate}
    de forma que:
    \begin{enumerate}
        \item $X(U) = V$.
        \item $X:U\to V$ es un homeomorfismo\footnote{Nos gustaría exigir que sea un difeomorfismo, pero $V$ no es necesariamente un espacio normado, con lo que no tenemos definido qué es ser una aplicación diferenciable $X:U\to V$.}.
        \item $(dX)_q:\mathbb{R}^2\to \mathbb{R}^3$ es inyectiva $\forall q\in U$.
    \end{enumerate}
\end{definicion}

\begin{observacion}
    Vemos que definimos anteriormente las curvas como aplicaciones, pero ahora las superficies las hemos dado como un subconjunto de $\mathbb{R}^3$, de forma que para cada punto $p\in S$ obtenemos una aplicación diferenciable $X:U\to \mathbb{R}^3$.\\

    \noindent
    La definición hipotética de que una superficie sea una aplicación diferenciable $X:U\to \mathbb{R}^3$ no nos gusta porque entonces la esfera no sería una superficice, sino que habría que cortar y pegar trozos, algo que no queremos hacer en geometría.\\

    \noindent
    Así, una superficie se obtiene juntando varios trozos de $\mathbb{R}^2$ y pegándolos juntos.
\end{observacion}

\begin{notacion}
    Usualmente notaremos a la aplicación $X:U\to \mathbb{R}^3$ que obtenemos en cada punto $p\in S$ por:
    \begin{equation*}
        X(u,v) = (x(u,v), y(u,v), z(u,v))
    \end{equation*}
    donde $x,y,z:U\to \mathbb{R}$ son aplicaciones diferenciables. Notaremos también:
    \begin{equation*}
        q = (u,v), \quad p = (x,y,z)
    \end{equation*}
\end{notacion}

\begin{observacion}
    Como $X:U\to V$ ha de ser un homeomorfismo, las superficies que estudiamos no pueden tener autointersecciones.
\end{observacion}

\begin{observacion}
    Otra forma de decir que cada aplcación $(dX)_q$ es inyectiva es decir que:
    \begin{equation*}
        \frac{\partial X}{\partial u}(q), \frac{\partial X}{\partial v}(q)
    \end{equation*}
    son linealmente independientes, que es tanto como el análogo a decir que tenemos una curva regular. Otra forma para ver esto es tanto como decir que no son linealmente dependientes, es decir, que:
    \begin{equation*}
        \frac{\partial X}{\partial u}(q) \land \frac{\partial X}{\partial v}(q) \neq 0 \quad \forall q\in U
    \end{equation*}
\end{observacion}

\begin{ejemplo}
    Ejemplos de superficies:
    \begin{enumerate}
        \item Los planos de $\mathbb{R}^3$ son superficies.

            Para verlo, sea $P\subset\mathbb{R}^3$ un plano:
            \begin{equation*}
                P = \{(x,y,z)\in \mathbb{R}^3 : ax + by + cz + d= 0\} \qquad (a,b,c)\neq (0,0,0), d\in \mathbb{R}
            \end{equation*}
            Como $(a,b,c)\neq (0,0,0)$ tendremos que al menos uno de ellos será no nulo. Supongamos que $c\neq 0$, con lo que:
            \begin{equation*}
                P = \{(x,y,z)\in \mathbb{R}^3 : z = Ax + By + C\} \qquad A,B,C\in \mathbb{R}
            \end{equation*}
            Sea $p\in P$, elegimos en vista de la definición:
            \begin{equation*}
                U = \mathbb{R}^2, \quad V = P, \quad X:U\to \mathbb{R}^3
            \end{equation*}
            dada por:
            \begin{equation*}
                X(u,v) = (u, v, Au + Bv + C)
            \end{equation*}
            observamos que es diferenciable. Para revisar las condiciones:
            \begin{itemize}
                \item $X(U) = X(\mathbb{R}^2) = \{(u,v,Au+Bv+C) \in \mathbb{R}^3 : (u,v)\in \mathbb{R}^2\} = P = V$.
                \item Veamos que $X:U=\mathbb{R}^2\to V=P$ es un homeomorfismo\footnote{No tenemos que demostrar ni la continuidad de $X$ ni su sobreyectividad, puesto que sabemos ya que $X(U) = V$ y que es un difeomorfismo.}, puesto que $X^{-1}:P\to \mathbb{R}^3$ dada por $(x,y,z)\mapsto (x,y)$ es su aplicación inversa, que es continua.
                \item $X_u(u,v) = (1,0,A)\quad \forall (u,v)\in U$ y que $X_v(u,v) = (0,1,B)\quad \forall (u,v)\in U$, que son linealmente independientes.
            \end{itemize}
        \item En realidad hemos probado algo más fuerte que lo de que un plano sea superficie, porque observamos que lo único que hemos usado de que $P$ sea un plano es que podíamos haber despejado $z$ en función de $x$ e $y$. Así, demostramos a continuación que el grafo de cualquier función real de dos variables diferenciable es una superficie:

            Sea $O\subset \mathbb{R}^2$ un abierto y $f:O\to \mathbb{R}$ una aplicación diferenciable, tenemos que:
            \begin{equation*}
                Gr f = \{(x,y,z)\in \mathbb{R}^3 : (x,y)\in O, z=f(x,y)\}
            \end{equation*}
            Fijado $p\in Gr f$, elegimos:
            \begin{equation*}
                U = O, \quad V = Gr f, \quad X:U\to \mathbb{R}^3
            \end{equation*}
            dada por:
            \begin{equation*}
                X(u,v) = (u,v,f(u,v))
            \end{equation*}
            que es una aplicación diferenciable. Revisamos las condiciones:
            \begin{itemize}
                \item $X(U) = X(O) = Gr f = V$.
                \item Veamos que $X:U\to V$ es un homeomorfismo, algo que se ha probado ya en anteriores cursos de Topología.
                \item Para ver que la diferencial sea inyectiva miramos las derivadas parciales:
                    \begin{equation*}
                        X_u(u,v) = (1,0,f_u(u,v)), \qquad X_v(u,v) = (0, 1, f_v(u,v))
                    \end{equation*}
                    Obtenemos vectores linealmente independientes.
            \end{itemize}
            Tanto el plano como el grafo tienen en común que $U,V,X$ no dependen del punto $p\in S$ escogido, por lo que estos dos objetos cumplirían aquella definición hipotética de una aplicación diferenciable de $U$ en $\mathbb{R}^3$.
        \item Sea $S$ una superficie y $O\subset S$ un abierto suyo, entonces $O$ es una superficie.

            Para verlo, sea $p\in O$, tenemos que $p\in S$ superficie, con lo que existen $U\subset \mathbb{R}^2$ abierto, $V$ entorno abierto de $p$ en $S$ y $X:U\to \mathbb{R}^3$ diferenciable cumpliendo 1, 2 y 3. Tomamos: $V_O = U\cap O$, que es un entorno abierto de $p$ en $V$, $U_O = X^{-1}(V\cap O)$, que por la condición 2 y ser $V\cap O$ un abierto de $V$ es un abierto de $\mathbb{R}^2$. Finalmente, $X_O = X\big|_{U_O}:U_O\to \mathbb{R}^3$ que sigue siendo diferenciable por el carácter local.

            Vemos ahora que: 
            \begin{equation*}
                X_O (U_O) = X\big|_{U_O}(X^{-1}(V\cap O)) = V\cap O = V_O
            \end{equation*}
            Así como que $X_O:U_O\to V_O$ es un homeomorfismo, puesto que es igual a:
            \begin{equation*}
                X\big|_{U_O}:X^{-1}(V\cap O) \to V\cap O
            \end{equation*}
            Finalmente, para $q\in U_O = X^{-1}(V\cap O)$ tenemos que:
            \begin{equation*}
                (dX_O)_q = (dX\big|_{X^{-1}(V\cap O)})_q = (dX)_q
            \end{equation*}
            Por lo que $(dX_O)_q:\mathbb{R}^2\to \mathbb{R}^3$ es inyectiva.

            Por tanto, acabamos de probar que $O$ es una superficie.
        \item Sea $S\subset\mathbb{R}^3$ no vacío, $S_i\subset S$ abierto de $S$ y superficie para cada $i \in I$ de forma que $S=\bigcup\limits_{i \in I}S_i$, entonces $S$ es una superficie.

            Para verlo, sea $p\in S=\bigcup\limits_{i \in I}S_i$, existe por tanto $i_0\in I$ de manera que $p\in S_{i_0}$ superficie, con lo que existen $U_0\subset \mathbb{R}^2$ abierto, $V_0$ entorno abierto de $p$ en $S_{i_0}$ y $X:U_0\to \mathbb{R}^3$ diferenciable de forma que verifican 1, 2 y 3.

            Elegimos $V=V_0\subset S$ entorno abierto de $p$ en $S$ por ser $V_0$ un entorno abierto de $p$ en $S_{i_0}\subset S$ siendo este último abierto; $U=U_0$ abierto de $\mathbb{R}^2$ y $X=X_0$ con lo que $X:U = U_0 \to \mathbb{R}^3$ es diferenciable. Falta ver que $U,V$ y $X$ verifican 1, 2 y 3. % // TODO: HACER

            % // TODO: Explicar que a partir de aqui se demustra que la union de superficies es superficie. Union de dos planos paralelos si pero union de dos planos que se cortan no por ser dos abiertos relativos.
        \item Sea $S\subset \mathbb{R}^3$ una superficie y $\phi:O\to O'$ un difeomorfismo con $O$ y $O'$ abiertos de $\mathbb{R}^3$ de forma que $S\subset O$, entonces $S' = \phi(S)$ es otra superficie en $\mathbb{R}^3$.

            Para ello, sea $p'\in S'$, existe un único $p\in S$ de manera que $\phi(p) = p'$. Como $S$ es superficie, asociados a $p$ tenemos $U,V$ y $X$ dados por la definición de superficie. Tomamos $U' = U$, $V' = \phi(V)$ y $X' = \phi\circ X$, tenemos que:
            \begin{itemize}
                \item $X'(U') = (\phi\circ X)(U) = \phi(V) = V'$.
                \item $X':U'\to V'$, es igual a $\phi\circ X:U\to \phi(V)$, que es composición de los homeomorfismos $X:U\to V$ y $\phi:V\to \phi(V)$, por lo que $X'$ es un homeomorfismo.
                \item Sea ahora $q'\in U'=U$, tenemos que:
                    \begin{equation*}
                        (dX')_{q'} = d(\phi\circ X)_{q'} = (d\phi)_{X(q')}\circ (dX)_{q'}
                    \end{equation*}
                    Siendo la última inyectiva por hipótesis y la primera biyectiva por ser una diferencial de un difeomorfismo, por lo que $(dX')_{q'}$ será también inyectiva.
            \end{itemize}
    \end{enumerate}
\end{ejemplo}

\begin{notacion}
    Si tenemos una superficie $S$, fijado $p\in S$ podemos obtener la $X$ que nos dice su definición, por lo que podemos escribir:
    \begin{equation*}
        X(u,v) = (x,y,z), \qquad X(q) = p
    \end{equation*}
    Llamaremos:
    \begin{itemize}
        \item A $(u,v)$ coordenadas de $p$.
        \item A $X$ parametrización (local), o a veces mapa o carta (de navegación).
        \item A $V$ entorno parametrizado.
        \item A $U$ dominio de la parametrización.
    \end{itemize}
\end{notacion}
